\documentclass[10pt,a4paper]{amsart}
\usepackage[margin=19mm]{geometry}
\usepackage{amsmath,amssymb,mathtools}
\usepackage[scr=rsfs,cal=euler]{mathalfa}
\usepackage{xfrac}
\usepackage[unicode,hidelinks,bookmarks=false]{hyperref}
\newcommand{\ie}{i.e.\ }
\newcommand{\cf}{cf.\ }
\newcommand{\resp}{respectively\ }
\newcommand{\Subsection}{\subsection}
\providecommand{\nobreakdash}{\nobreak\hbox{-}}
\setlength{\emergencystretch}{2em}
\multlinegap0pt
\usepackage{amssymb}
\usepackage{algorithm}\usepackage{algpseudocode}
\newtheorem{lemma}{Lemma}
\newtheorem{theorem}{Theorem}
\newcommand{\ad}{\ensuremath{\mathrm{ad}}}
\newcommand{\e}{\ensuremath{\mathrm{e}}}
\title{Error bounds for the truncated Baker--Campbell--Hausdorff and Zassenhaus formulas in unitary problems}
\author{A. Arnal, F. Casas, J.L. Ruiz-Benito}
\date{Source: arXiv:2607.07692v1 (CC BY 4.0)}
\begin{document}
\maketitle

\noindent\emph{Source and licence.} Error bounds for the truncated Baker--Campbell--Hausdorff and Zassenhaus formulas in unitary problems. Version arXiv:2607.07692v1 (CC BY 4.0), distributed by arXiv under the Creative Commons Attribution 4.0 International licence, http://creativecommons.org/licenses/by/4.0/, as recorded in the arXiv licence metadata for this exact version and shown on the abstract page https://arxiv.org/abs/2607.07692v1. Full licence notice: \texttt{licenses/arxiv-2607.07692-CC-BY-4.0.txt}.\par\smallskip

\noindent\textit{Editorial scope.} Counted unit: the article's theorem th:general\_bound (source0.tex lines 389--401). The article's complete original proof of it is delivered with it (source0.tex lines 403--435, one \texttt{proof} environment of 33 lines). No mathematics is written, repaired or re-derived: the statement and the proof are the article's own lines, in the article's own notation, and no retained line is substituted or re-flowed.  Retained uncounted as the source-local context the proof needs: lines 230--240; lines 265--280; lines 357--359.  Nothing was omitted from any retained span, so the package carries no omission record.  No retained line contains a citation, so the wrapper carries no bibliography and no bibliography entry.  No retained line contains a figure environment, an image-inclusion command or drawing code, so no image is delivered and the package declares no asset dependency.  Editorial formatting only, in addition: an AMS \texttt{amsart} preamble that restates the article's own declarations and macros used by the retained lines, namely the environments \texttt{lemma}, \texttt{theorem} on separate counters, together with the standard packages those lines need.  The retained environments print in this document as Lemma 1, Theorem 1 in order of appearance, whereas the article prints the same items with the numbering of its own full document.  The complete native source of the article as distributed in the arXiv e-print for this exact version is bundled unchanged as \texttt{sources/204732/source0.tex}.
\begin{lemma} \label{lemma_estimate}
Let $F_1(t)$, $F_2(t)$ be two unitary operators verifying the initial value problems
\begin{equation} \label{ivp_2op}
   \frac{d F_i}{dt} = M_i(t) F_i, \qquad F_i(0) = I, \qquad\qquad i=1,2,
\end{equation}
with $M_i(t)$ skew-adjoint operators. Then,
\[
\|F_1(t) - F_2(t)\| \leq 
\int_0^t \| M_1(s) - M_2(s)\| \, ds.
\]
\end{lemma}  

\begin{lemma} \label{l:lemma2}
    Let $X, Y$ be two linear operators such that $\e^{s X}$ is well defined for all $s \in \mathbb R$. Then, for any $n\geq 1$, 
        \[
        \e^{s \, \ad_X} Y = \sum_{k=0}^{n-1} \frac{s^k}{k!} \, \ad_X^{k} Y + s^n G_n\left(s, X, Y\right),
    \]
    where
  \begin{equation}
  \label{eq:Gn}
G_n(s, X,Y) = \int_0^1 \frac{\sigma^{n-1}}{(n-1)!} \, \e^{s \, (1-\sigma)\, \ad_X}\, \ad_X^n Y\, d\sigma.
\end{equation}
If in addition $X$ is skew-adjoint, then
\begin{equation} \label{norm_G}
  \|G_n(s,X,Y)\| \le \int_0^1 \frac{\sigma^{n-1}}{(n-1)!} \|\ad_X^n Y\| d\sigma \le \frac{1}{n!} \|\ad_X^n Y\|.
\end{equation}

\end{lemma}

 \begin{equation} \label{M1M2_trunc}
   M_1(t) = A + \e^{t \, \ad_{A}} B, \qquad\qquad  M_2(t) = \int_0^1 \e^{x \, \ad_{Z^{[m]}}} \big(\dot{Z}^{[m]}\big) dx,
 \end{equation}

\begin{theorem}\label{th:general_bound} 
Let $A$ and $B$ be skew-adjoint operators, and $Z^{[m]}(tA, tB)$ denote the sum of the first $m \ge 2$ 
terms of the Baker--Campbell--Hausdorff series and $t > 0$. Then
\[
   \left\| \e^{t A} \, \e^{t B} - \e^{Z^{[m]}(tA, tB) } \right\| \le \int_0^t \|  \mathcal{R}^{[m]}(s) \| ds + \frac{t^{m+1}}{(m+1)!} \|\ad_A^m B\| + 
   \frac{1}{m!} \int_0^t \|\ad_{Z^{[m]}}^{m-1} \dot{Z}^{[m]} \| ds
\]
where 
\begin{equation} \label{def_R}
    \mathcal{R}^{[m]}(s) = A + \sum_{k=0}^{m-1} \frac{s^k}{k!} \ad_A^k B - \sum_{k=0}^{m-2} \frac{1}{(k+1)!} \, \ad_{Z^{[m]}}^k \dot{Z}^{[m]} ,
\end{equation}
and $ \dot{Z}^{[m]}$ denotes the derivative of $Z^{[m]}(sA,sB)$ with respect to $s$.
\end{theorem}   

\begin{proof}
Application of Lemma \ref{l:lemma2} to \eqref{M1M2_trunc} gives
\[
 \begin{aligned}
  & M_1(s) = A + \sum_{k=0}^{m-1} \frac{s^k}{k!} \ad_A^k B + s^m G_m(s,A,B) \\
  & M_2(s) =  \sum_{k=0}^{m-2} \frac{1}{(k+1)!} \ad_{Z^{[m]}}^k \dot{Z}^{[m]} + \int_0^1 x^{m-1} G_{m-1}(x,Z^{[m]},  \dot{Z}^{[m]}) dx,
 \end{aligned}
\]
so that
\[
 \|M_1(s) - M_2(s)\| \le \|  \mathcal{R}^{[m]}(s) \| + s^m \|G_m(s,A,B)\| +  \int_0^1 x^{m-1} \| G_{m-1}(x,Z^{[m]},  \dot{Z}^{[m]})\| dx.
\]
Notice that $\mathcal{R}^{[m]}(s)$ defined by \eqref{def_R} verifies 
\[
   \mathcal{R}^{[m]}(s) = \mathcal{O}(s^m),
\]
since all terms up to $s^{m-1}$ vanish due to the way the successive terms $\Phi_1, \ldots, \Phi_m$ in the BCH series are obtained.  In fact, 
the term in $s^m$ comes from
$\sum_{k=1}^{m-2} \frac{1}{(k+1)!} \, \ad_{Z^{[m]}}^k \dot{Z}^{[m]}$. 


Now, from \eqref{norm_G}, we have
  $\|G_m(s,A,B)\| \le  \frac{1}{m!} \|\ad_A^m B\|$
and
\begin{equation} \label{norm_G2}
\begin{aligned}
   & \int_0^1 x^{m-1} \| G_{m-1}(x,Z^{[m]},  \dot{Z}^{[m]})\| dx   \le  \int_0^1 \frac{x^{m-1}}{(m-1)!} 
    \| \ad_{Z^{[m]}}^{m-1} \dot{Z}^{[m]} \| \, dx   \\
  & \qquad\quad = \frac{1}{m!} \| \ad_{Z^{[m]}}^{m-1} \dot{Z}^{[m]} \|.
\end{aligned}    
\end{equation}  
Lemma \ref{lemma_estimate} leads finally to the desired result. 
\end{proof}   

\end{document}
